\section{Conclusions}

In this laboratory work for the course \enquote{Numerical Methods} I implemented the basic
methods of numerical linear algebra and applied them to variant 28.

Direct methods give the solution in a fixed number of operations. The LU decomposition with
partial pivoting solves a general system in $O(n^3)$, and the same factorization gives the
determinant and the inverse matrix. The Thomas algorithm uses the structure of a tridiagonal
matrix and solves the system in $O(n)$; diagonal dominance with non-zero off-diagonal elements
guarantees its stability.

Iterative methods trade exactness for simplicity of every step. For the diagonally dominant
system of the variant the Seidel method needed about half as many iterations as the simple
iteration method, in agreement with the ratio of the spectral radii of their iteration matrices,
while the a priori estimate turned out to be about 7 times too pessimistic.

The rotation method finds the complete eigensystem of a symmetric matrix: the error decreases
monotonically and at the end converges quadratically. The QR algorithm with Householder
reflections works for arbitrary matrices and finds complex eigenvalues as $2 \times 2$ blocks;
its rate depends on the ratio of the moduli of the eigenvalues.

All results were verified by unit tests on random matrices (residuals, orthogonality, trace and
determinant) and by end-to-end tests of the programs.

\pagebreak
